% !TeX spellcheck = de_DE

Die Qualitätsbewertung ist das Herzstück des gesamten Ansatzes.
Hierfür werden die vom Metrics Service bereitgestellten Kennzahlen abgerufen und mit zuvor definierten Zielwerten verglichen.

Für jede Metrik wird geprüft, ob die entsprechenden Anforderungen erfüllt werden.
Das Ergebnis dieser Prüfung dient als Indikator für die aktuelle Qualität des Textes.
Gleichzeitig werden erkannte Abweichungen als Feedback für den Re-Prompt gespeichert und aufbereitet.

Die konkrete Auswahl der verwendeten Metriken ist Bestandteil der Untersuchung und wird im Rahmen der Evaluation analysiert.
Der Prototyp ist daher so konzipiert, dass unterschiedliche Metriken flexibel integriert und ausgetauscht werden können.